\section{Approach}

\subsection{One rule}

The pipeline is built on a single constraint: \textbf{every stage must either
verify its input against a stated precondition, or decline it}. Declining is
not an empty result. It is an output, carrying a machine-readable reason, and
it is counted in the same ledger as an admitted record.

This is a stronger commitment than it sounds. No stage may fall through,
guess, or interpolate. A symbol that does not resolve is not approximately
matched. A record whose numbering disagrees with its reference is not applied
anyway. A source that a third party cannot re-query is not consulted at all.

The input is a rare-disease name. The output is a set of (disease, target,
drug-programme) rows. Each row carries a verdict in $\{\textsc{gap},
\textsc{not\_a\_gap}, \textsc{abstain}\}$, the identifier of the rule that
produced it, and a receipt: the content digests of every record the verdict
consumed.

\subsection{Six types of evidence}

We collect six kinds of evidence, each answering exactly one question
(Table~\ref{tab:evidence}). The separation is deliberate. A source is
consulted for the one thing it is authoritative about and for nothing else, so
that a failure in one lane cannot silently substitute for another.

\begin{table}[h]
\centering
\small
\caption{The six evidence types. Each source answers one question, and none is
consulted outside its competence.}
\label{tab:evidence}
\begin{tabular}{lll}
\toprule
Evidence type & Question it answers & Admitted as \\
\midrule
Target association  & which genes are implicated in this disease?       & target set \\
Target tractability & is there existing chemistry against them?         & druggability \\
Clinical variants   & do patients carry damaging variants here?         & genetic support \\
Canonical sequence  & what is the reference protein?                    & substitution frame \\
Registered trials   & has this pairing been tried, and how did it end?  & prior attempt \\
Regulatory status   & is the drug approved, and for what?               & repurposing cost \\
Cell morphology     & does any compound move cells toward reference?    & phenotype rank \\
\bottomrule
\end{tabular}
\end{table}

Two details matter for the evaluation that follows. First, trials retain
\emph{status and phase}, not merely existence. A terminated Phase~3 is evidence
against a pairing. A pipeline that records only that a trial exists converts a
past failure into an apparent opportunity.

Second, the canonical sequence is treated as evidence in its own right rather
than as a lookup. It is the frame against which every variant record is
checked, and the first guard we ablate is precisely a check against it.

\subsection{The provenance framework}

Every response is hashed the moment it arrives, and its query URL is recorded
beside the digest. Records carry one of two custody levels.

\begin{itemize}
\item \textbf{Recomputed.} We hold the bytes, and the digest can be re-derived
from them locally.
\item \textbf{Committed.} We hold a digest captured at origin but not the
bytes. This is the honest level for a large third-party artifact we do not
redistribute.
\end{itemize}

The distinction is rendered rather than hidden, because the two levels support
different claims. Every record, at either level, is content-addressed and
committed under a single Merkle root, using the audit-tree construction of
Certificate Transparency \citep{rfc6962}. Verification re-derives the root in
the reader's browser from the record bytes. A single-byte edit moves the root,
and the first divergent node is named.

Admissibility follows from this. A source whose responses a third party cannot
re-obtain cannot be given a meaningful digest, and is therefore inadmissible.
One candidate source is excluded on exactly this ground. We report it as
excluded rather than omitting it quietly.

\subsection{Dispositions: every record lands somewhere}

The provenance framework is half of what the evaluation needs. The other half
is a \emph{closed disposition vocabulary}. Each stage assigns every record it
sees to exactly one labelled outcome, and the labels are exhaustive.

The variant lane admits a record, excludes it as a multi-gene copy-number
event, or reports it unavailable from the source. The sequence lane emits a
substitution or a truncation, or declines with one of six named reasons
(Table~\ref{tab:decline}). The decision layer assigns one of five verdicts
through nine first-match rules. No stage has an unlabelled path.

Each stage then asserts its own arithmetic: records in must equal records out
plus records declined. For the corpus used here, acquisition reconciles as
\CorpusFetched{} fetched $=$ \CorpusKept{} admitted $+$ \CorpusExcluded{}
excluded $+$ \CorpusUnavailable{} unavailable. The sequence lane reconciles as
\GOneProcessed{} in $=$ \GOneEmits{} emitted $+$ \SeqDeclined{} declined.

\subsection{Why provenance is a precondition, not an adjunct}

Custody is usually motivated by trust: a reader who does not trust the authors
can check the evidence. That holds here, but it is not why the framework is
load-bearing in this paper.

\textbf{The ablation is measurable only because the pipeline is fully
accounted.} To report what an unguarded pipeline \emph{would have emitted}, we
must know the complete set of records it would have seen and the disposition
each one actually received. A pipeline that drops records on an unlabelled
path cannot support that counterfactual. The denominator is unknown, and no
amount of care in reading the code recovers it.

This is not hypothetical. Section~\ref{sec:audit} reports that our own
acquisition stage was discarding 100 records through an unlabelled branch. The
per-gene counts looked entirely normal. The reconciliation assertion exposed
the gap, and review had not. Until it was closed, every rate in this paper had
an unknown denominator.

\subsection{Decision}

Nine first-match rules assign the verdict. They cover unresolved symbols,
absent targets, absent programmes, failed lookups, prior trials, and three
grades of untried pairing. The rules are plain code and are read in order. No
model output may set a verdict, resolve a symbol, or choose a threshold. A
language model is used only to phrase rows whose verdict is already fixed.

Identity resolution deserves its own note, because it is where silent errors
originate. Symbols resolve by exact match or an explicit alias table, and by
nothing else. Collisions are real. One symbol in our corpus denotes two
unrelated proteins, and a similarity match would return a valid gene, succeed
at every downstream step, and answer a question about the wrong molecule.

\subsection{The guards under test}
\label{sec:guards}

Two guards carry the measurements in Section~\ref{sec:eval}.

\textbf{G1, residue verification.} A variant record names a wild-type residue
and a position. Before applying the substitution to the canonical sequence,
the pipeline checks that the named residue is what the canonical sequence
actually holds at that position. If it is not, the record is numbered against
a different transcript, and the row declines.

\textbf{G2, copy-number exclusion.} A variant record is admitted as evidence
about a gene only if it is gene-specific. Records typed as copy-number events,
or spanning more than a small number of genes, are excluded. A deletion
covering hundreds of genes is not evidence about any one of them.
